\clearpage\fancyhead[L]{\small COMET.}\hypertarget{native-348-lower}{}\label{native-348-lower}
\begin{center}\textbf{COMET.}\par\vspace{4pt}\hrule\vspace{4pt}\fontsize{10}{12}\selectfont Monatliche Correspondenz. Vol. XXIV, pp. 595\,..\,598, December 1811.\par\vspace{4pt}\hrule\vspace{4pt}\end{center}

\par\indent In Göttingen, Messrs. \textsc{Gauss} and \textsc{Harding} found the comet on the very evening my letter arrived there, 9 December. \textsc{Gauss} made the following observations with the ring micrometer.

\begin{center}
\fontsize{10}{12}\selectfont\setlength{\tabcolsep}{4pt}\renewcommand{\arraystretch}{1.15}\begin{tabular}{l|c|c|c}
 & Mean time & Comet R.A. & S. decl. \\
 & in Göttingen & & \\ \hline
Dec. 9 & $10^{\mathrm u}\;\;6^{\mathrm m}\;52^{\mathrm s}$ & $63^\circ\;49'\;41''4$ & $10^\circ\;21'\;55''5$ \\
11 & $10\;\;34\;\;\;\;1$ & $63\;\;33\;\;18.0$ & $8\;\;39\;\;46.4$ \\
12 & $\;\;8\;\;\;\;5\;\;52$ & $63\;\;26\;\;25.3$ & $7\;\;54\;\;25.9$ \\
\end{tabular}
\end{center}

\SourceLeaf{349}
\par\indent “Here, wrote \textsc{Gauss} to us from Göttingen on 15 Dec. 1811, are also my provisional elements. They can of course be regarded only as an approximation, but they do show that the comet is a \emph{new} one and will grow no brighter.

\Needspace{16\baselineskip}
\begin{center}\fontsize{10}{12}\selectfont\setlength{\tabcolsep}{4pt}\renewcommand{\arraystretch}{1.15}\begin{tabular}{l}
Time of perihelion, 1811 Nov. 12,6225, Göttingen meridian\\
\end{tabular}\end{center}
\begin{center}\fontsize{10}{12}\selectfont\setlength{\tabcolsep}{3pt}\renewcommand{\arraystretch}{1.15}\begin{tabular}{llr}
Longitude of perihelion & .\;.\;.\;.\;.\;.\;. & $48^\circ\;30'\;20''$\\
Log. of the perihelion distance & .\;.\;.\;.\;. & $0.20160$\\
Ascending node & .\;.\;.\;.\;.\;. & $92^\circ\;46'\;59''$\\
Inclination of the orbit & .\;.\;.\;.\;.\;.\;. & $31\;37\;55$\\
\end{tabular}\end{center}
\noindent Direct motion


\par\indent “The comet may remain visible for a \emph{very} long time, if its brightness allows. For 31 January 1812, $10^{\mathrm u}$, I find R.A. $67\frac{1}{2}^\circ$, decl. $25\frac{1}{2}^\circ$ north, brightness 0.178.

\begin{center}
\begin{tabular}{l}
Brightness 1811 November 18 $\;=\;0.713$ \\
\qquad --- \qquad --- December 12 $\;=\;0.698$ \\
\end{tabular}
\end{center}


\Needspace{20\baselineskip}\par\indent Mr. \textsc{Gerling} compared all the observations that had become known by then with these elements and found the following results.

\begin{center}
\fontsize{10}{12}\selectfont\setlength{\tabcolsep}{6pt}\renewcommand{\arraystretch}{1.15}\begin{tabular}{l|c|c|l}
 & \multicolumn{2}{c|}{Deviation} & \\
 & in R.A. & in decl. & Observer \\ 
\hline
Nov. 18 & $+32''$ & $-101''$ & v. Zach \\
19 & $+80$ & $+46$ & v. Zach \\
20 & $+23$ & $+4$ & v. Zach \\
21 & $+3$ & $+15$ & v. Zach \\
Dec. 8 & $+19$ & .\;.\;.\;. & v. Lindenau \\
9 & $+11$ & $+12$ & v. Lindenau \\
9 & $+8$ & $-55$ & Olbers \\
9 & $+1$ & $+1$ & Gauss \\
11 & $-10$ & $-16$ & Gauss \\
12 & $+1$ & $+7$ & Gauss \\
\end{tabular}
\end{center}


\par\indent To facilitate future observations of the comet, Mr. \textsc{Nicolai} used the elements above to calculate the following ephemeris.

\begin{center}
\fontsize{10}{12}\selectfont\setlength{\tabcolsep}{4pt}\renewcommand{\arraystretch}{1.15}\begin{tabular}{l|c|c|c|c}
Midnight & R.A. & Declination & \shortstack{Log. dist.\\$a\;\EarthSymbol$} & \shortstack{Brightness} \\ \hline
1811 December 21 & $62^\circ\;36'$ & $0^\circ\;7$ S. & 9.8879 & 0.596 \\
29 & $62\;\;25$ & $6\;\;20$ N. & 9.9185 & 0.496 \\
1812 January 6 & $62\;\;49$ & $12\;\;11$ --- & 9.9562 & 0.397 \\
14 & $63\;\;47$ & $17\;\;15$ --- & 9.9984 & 0.310 \\
22 & $65\;\;17$ & $21\;\;33$ --- & 0.0427 & 0.239 \\
30 & $67\;\;17$ & $25\;\;\;\,8$ --- & 0.0872 & 0.183 \\
February 7 & $69\;\;41$ & $28\;\;\;\,7$ --- & 0.1310 & 0.141 \\
\end{tabular}
\end{center}


\par\indent According to these figures, the comet’s brightness is already decreasing quite noticeably again. Since moonlight will now hinder its visibility for a while, it still seems doubtful whether another series of good observations of this comet can be obtained. Yet the fine agreement with all the observations shows that the elements above already determine its course to a very close approximation, so our descendants will be able to recognize it without difficulty if it reappears someday.
